% TOP_PROBLEM: {"id": 3008, "title": "Kirby Problem 5.1", "category_id": 7, "category": {"id": 7, "name": "topology", "display_name": "Topology", "description": "Properties preserved under continuous deformations.", "slug": "topology", "order_index": 7, "created_at": "2026-07-31T15:26:25.671Z"}, "status": "open", "problem_number": "KP-5.1", "set_id": 11}
% FIRST_POSTED: 2026-10-01
% ENTRY_KIND: statement_only
% UnsolvedMath ID 3008; source catalogue code KP-5.1
% !TeX program = lualatex
% Definition and sources only; this file is not an LLM solution attempt.
\documentclass[11pt,a4paper]{article}
\usepackage[margin=25mm]{geometry}
\usepackage{fontspec}
\setmainfont{lmroman10-regular.otf}[BoldFont=lmroman10-bold.otf,
 ItalicFont=lmroman10-italic.otf,BoldItalicFont=lmroman10-bolditalic.otf]
\usepackage{amsmath,amssymb,mathtools}
\usepackage{unicode-math}
\setmathfont{latinmodern-math.otf}
\AtBeginDocument{\let\setminus\smallsetminus}
\usepackage{xurl,hyperref}
\hypersetup{colorlinks=true,urlcolor=blue}
\setcounter{secnumdepth}{0}
\setlength{\parindent}{0pt}
\setlength{\parskip}{5pt}
\setlength{\emergencystretch}{3em}
\begin{document}
\section*{Kirby Problem 5.1}\label{3008}
{\small UnsolvedMath ID 3008 · KP-5.1 · Topology · Kirby's Problems in Low-Dimensional Topology}
\subsection{Definitions and mathematical statement}
A subset $C\subset\mathbb R^2$ is cellular if it can be written
as a nested intersection
\[
 C=\bigcap_{n=1}^{\infty}D_n,\qquad
 D_{n+1}\subset\operatorname{int}(D_n),
\]
where each $D_n$ is an embedded closed topological $2$-disk in
the plane. An embedded disk is the image of the standard closed
disk under a topological embedding; no smoothness of its boundary
is required. In the plane, this describes the nonempty compact
connected sets whose complement is connected.

A space has the fixed point property if every continuous map
$f:C\to C$ has a point $x\in C$ with $f(x)=x$.
Does every planar cellular set have this property? The map is
an arbitrary continuous self-map, with no injectivity, surjectivity,
or extension-to-the-plane hypothesis.

The set $C$ need not be a disk, a manifold, locally connected,
or locally contractible. It may have complicated limiting geometry
even though each approximating $D_n$ is a disk. Compactness
ensures the nested intersection is nonempty, but the requested
fixed point must lie in $C$ and be fixed by the given map on
$C$ itself.

The assertion is universal in both the cellular set and the
continuous self-map. A counterexample would therefore consist of
one planar cellular continuum with a fixed-point-free continuous
self-map. The question is specifically two-dimensional; a cellular
set embedded in a higher-dimensional Euclidean space belongs to
a different problem.
\subsection{Short English statement}
Does every continuous self-map of every cellular subset of the plane have a fixed point?
\subsection{Sources}
\begin{itemize}
\item[S1] ULAM.AI and the UnsolvedMath Contributors, supplied problems.json, record 3008 (KP-5.1), Kirby's Problems in Low-Dimensional Topology. Catalogue identity and source transcription; CC BY 4.0.
\url{https://huggingface.co/datasets/ulamai/UnsolvedMath}.
\item[S2] K3: A New Problem List in Low-Dimensional Topology, Problem 5.1. Expanded formulation of the supplied catalogue statement.
\url{https://math.berkeley.edu/sites/default/files/surv-295-ruberman-watermarked-author-pdf.pdf}.
\end{itemize}
\end{document}
